% !TEX root = ../../../main.tex
% Block card: preprocessing.trace_shape -- "Trace shape" (Part I, the blocks chapter).
\cardfield{What it does}
Turns every window of a pool into its \emph{shape}: the trace stretched or squeezed to the
same number of points, with its size and level taken out.

\cardfield{When to use it}
Between \blockref{preprocessing.window_pool} and \blockref{catalogue.shape_cluster} (or
\blockref{catalogue.shape_families}), when you want a one-minute and a thirty-minute window
compared by the form of their trace alone. It is the same preparation the Library uses for
its families, so a distance here means what it means there.

\begin{plainterms}[What ``shape'' means once size is taken out]
Imagine tracing the outline of a hill on paper, then photocopying it to fit one page,
whatever its real width, and adjusting the copier so its height always fills the page. Two
hills of different size but the same profile now look the same. That is what this block does
to a window: it resamples it to 256 points (the width), subtracts its average and divides by
its spread (the height). What is left is only the shape. The cost: a tiny wiggle and a huge
drop with the same profile become identical, so windows too small to be more than noise are
left out first.
\end{plainterms}

\cardfield{What goes in, what comes out}
In go the windows of a pool. Out come the same windows, minus the ones under the noise floor,
each with its raw size (peak to peak, in millivolts) kept beside it. The shapes themselves are
written to a file on disk, which the next block reads.

\cardfield{Settings that matter}
\emph{align} says what each window is lined up on. \texttt{centre} (the default) moves each
window so its largest swing, up or down, sits in the middle; \texttt{peak} lines it up on its
highest point (the top of a spike); \texttt{trough} on its lowest point (the bottom of a drop);
\texttt{grid} keeps the windows where the pool cut them. Every re-cut stays inside its own
train or test stretch. Pick \texttt{peak} for spikes and \texttt{trough} for drops when a
window can hold both: \texttt{centre} follows whichever is larger, so a big spike in front of a
small drop puts the spike in the middle of one window and the drop in the middle of the next.
\emph{anchor\_frac} says where in the window the chosen point lands: 0.5, the middle, by
default; 0.33 puts it a third of the way in, which leaves room for a slow recovery after a
drop. It does nothing on \texttt{grid}. Change either and the shapes change, so the clusters
and families built on them do too. \emph{detrend}: \texttt{linear} removes a straight-line drift first, so a
slow slope is not mistaken for the shape. \emph{noise\_floor} leaves out windows smaller than
the dataset's noise floor (Settings, Datasets; 0.1\,mV when none is set).

\cardfield{Watch out for}
A recording with no declared unit has no size in millivolts: its windows are kept, never
tested against the floor, and counted as ``unmeasured''. Any re-cut (\texttt{centre}, \texttt{peak} or \texttt{trough}) can
make two windows land on the same event; one of them is dropped and counted. \texttt{peak} and
\texttt{trough} take the highest or lowest point whatever its size: a small spike before a big
drop is still what \texttt{peak} lines up on.
